\documentclass[crop=false]{standalone}
\usepackage{15150toolbox}
\usepackage{comment}
\usepackage{needspace}

% Toggle: set \soltrue to render solutions inline (blue boxes).
\newif\ifsol
\soltrue   % comment this line for the student version

% Solution environment.
\ifsol
  \newenvironment{solution}
    {\par\color{solution_color}\begin{framed}\textbf{Solution.}\par}
    {\end{framed}}
\else
  \excludecomment{solution}
\fi

% Lecture-slide macros not in 15150toolbox.
\newcommand{\term}[1]{\textbf{\textit{#1}}}
\newcommand{\thmBox}[2]{\par\noindent\begin{mdframed}[backgroundcolor=blue!8,hidealllines=true]\textbf{Theorem#1.} #2\end{mdframed}}
\newcommand{\bcBox}[1]{\par\noindent\textbf{Base case.} #1\par}
\newcommand{\ihBox}[2]{\par\noindent\textbf{Inductive hypothesis#1.} #2\par}
\newcommand{\isBox}[2]{\par\noindent\textbf{Inductive step#1.} #2\par}

% Problem header.
\newcommand{\quizproblem}[3]{%
  \needspace{5\baselineskip}%
  \stepcounter{problem}%
  \ifsol
    \subsection*{Problem \arabic{problem}. #1 \hfill \timing{#2} \quad (#3 pts)}%
  \else
    \subsection*{Problem \arabic{problem}. #1 \hfill (#3 pts)}%
  \fi
}

\newcommand{\subpts}[1]{\textbf{(#1 pts)}}

\newcommand{\answerspace}[1]{\ifsol\else\vspace{#1}\par\fi}

\printout{Quiz 5}

\begin{document}

\begin{center}
  {\Large \textbf{15-150 --- Quiz 5}} \\
  \vspace{4pt}
  Topics: \textit{Functors}; \textit{Red-Black
  Trees}; \textit{Sequences}; \textit{Lazy Programming} \\
  \vspace{4pt}
  \textit{60 minutes.}
\end{center}

\vspace{1em}

\noindent\textbf{Name:} \hrulefill \hspace{2em} \textbf{Andrew ID:} \hrulefill

\vspace{1em}

% =============================================================================
% PROBLEM 1 --- Functors (write an ORD instance + apply MkDict)
% =============================================================================

\quizproblem{Keys to the Kingdom}{7}{6}

Recall the type class of ordered types, and the functor \code{MkDict}
that turns any ordered key type into a dictionary:

\begin{codeblock}
  signature ORD =
    sig
      type t
      val compare : t * t -> order
    end

  functor MkDict (Key : ORD) :> POLY_DICT where type Key.t = Key.t =
    struct ... end
\end{codeblock}

Given a \code{structure K : ORD}, the application \code{MkDict (K)}
produces a dictionary whose keys are values of type \code{K.t}, kept
sorted according to \code{K.compare}. The SML basis provides
\code{Int.compare : int * int -> order}.

\textbf{(a)} \subpts{4} We want a dictionary keyed by \code{int}, but
sorted in \emph{descending} order (largest key first). Write a structure
\code{RevOrd} that ascribes to \code{ORD} with \code{type t = int} and a
\code{compare} that reverses the usual integer ordering, then use
\code{MkDict} to define \code{RevDict}. You may use \code{Int.compare}.

\begin{codeblock}
  structure RevOrd =







  structure RevDict =

\end{codeblock}

\answerspace{1.5cm}

\begin{solution}
\begin{codeblock}
  structure RevOrd : ORD =
    struct
      type t = int
      fun compare (x, y) = Int.compare (y, x)
    end

  structure RevDict = MkDict (RevOrd)
\end{codeblock}
Swapping the arguments to \code{Int.compare} reverses the order: now
\code{compare (x, y)} is \code{LESS} when \code{x > y}, so larger keys
sort first. Swapping \code{LESS}/\code{GREATER} in an explicit
\code{case} works too.

\textbf{Rubric:} 3, reversed \code{compare} (\code{Int.compare (y, x)} or
swapped cases; unreversed = 0); 1, \code{MkDict (RevOrd)}.
\end{solution}

\textbf{(b)} \subpts{2} Your \code{RevOrd} must be ascribed
\emph{transparently} (\code{:}), not opaquely (\code{:>}). In one
sentence, say what would go wrong for a client of \code{RevDict} if you
ascribed \code{RevOrd} opaquely instead.

\answerspace{3cm}

\begin{solution}
Opaque ascription hides \code{RevOrd.t = int}, making the key type
abstract. A client then has no values of that type to work with:
\code{RevDict.insert (5, v)} fails to typecheck, since \code{int} is not
known to match the key type.

\textbf{Rubric:} 2, opaque hides \code{RevOrd.t = int}, so a client
can't supply a real \code{int} key --- insert/lookup won't typecheck.
\end{solution}

\newpage

% =============================================================================
% PROBLEM 2 --- Red-Black Trees (definitional spot-the-bug on invariants)
% =============================================================================

\quizproblem{Seeing Red}{5}{4}

A \term{red-black tree} is a BST whose nodes are colored \code{Red} or
\code{Black}:

\begin{codeblock}
  datatype 'a rbtree =
      Empty
    | Red   of 'a rbtree * (Key.t * 'a) * 'a rbtree
    | Black of 'a rbtree * (Key.t * 'a) * 'a rbtree
\end{codeblock}

\textbf{(a)} \subpts{2} A student states the three red-black tree
invariants as follows. \textbf{Exactly one} of them is stated
incorrectly. Circle the number of the wrong one, and rewrite it so it
is correct.

\begin{enumerate}
  \item The tree is a BST: its inorder traversal is sorted with respect
    to its keys.
  \item The children of a \code{Red} node must be \code{Black}.
  \item Every path from the root to an \code{Empty} leaf passes through
    the same number of nodes.
\end{enumerate}

\answerspace{4.5cm}

\begin{solution}
\textbf{Invariant 3 is wrong.} It should count only \code{Black} nodes:
\emph{every path from the root to an \code{Empty} leaf passes through the
same number of \code{Black} nodes} (the \term{black height}). Counting
all nodes would force perfect height balance; red nodes are what let path
lengths differ. Invariants 1 and 2 are correct.

\textbf{Rubric:} 1, picks 3; 1, correction restricts the count to
\code{Black} nodes.
\end{solution}

\textbf{(b)} \subpts{2} When \code{insert} adds a new node, it always
colors that node \code{Red}, never \code{Black}. In one sentence, say
which invariant coloring the new node \code{Black} would immediately
risk breaking, and why.

\answerspace{3.5cm}

\begin{solution}
A new black node adds to the black count of its own root-to-leaf path
but no other, breaking the \textbf{black height invariant}. A red node
adds to no path's black count, so it preserves black height (it may
cause a red-red violation, which rebalancing fixes).

\textbf{Rubric:} 2, black height invariant --- a new black node
lengthens the black count on its path only.
\end{solution}

\newpage

% =============================================================================
% PROBLEM 3 --- Sequences (work/span recognition + associativity)
% =============================================================================

\quizproblem{The Cost of Doing Business}{7}{6}

Recall that \term{sequences} support the same operations as lists but
with a different \emph{cost model} that exposes parallelism. The
combinator \code{Seq.reduce g z S} combines all the elements of \code{S}
using \code{g}. Consider the function \code{maxS}, which finds the
largest element of a nonempty \code{int seq} (\code{NEG\_INF} is a value
smaller than every real element, and \code{Int.max} is constant work and
span):

\begin{codeblock}
  fun maxS (S : int seq) : int = Seq.reduce Int.max NEG_INF S
\end{codeblock}

Let \code{n = Seq.length S}.

\textbf{(a)} \subpts{4} Give a tight big-$O$ bound for the \textbf{work}
and the \textbf{span} of \code{maxS S}, and explain in one sentence
\emph{why} the span is what it is.

\vspace{0.3em}
\hspace{2em}Work: \rule{5cm}{0.4pt} \qquad Span: \rule{5cm}{0.4pt}

\answerspace{4cm}

\begin{solution}
\textbf{Work:} $O(n)$. \qquad \textbf{Span:} $O(\log n)$.

\textbf{Why the span is $O(\log n)$:} \code{reduce} combines pairwise up
a balanced tree, not one-at-a-time left to right like a list \code{fold}.
Each of the $\log n$ levels does its \code{Int.max}es in parallel, so the
longest dependency chain has length $O(\log n)$. Work stays $O(n)$: all
$n - 1$ combines still happen, parallelism just shortens the span.

\textbf{Rubric:} 1, work $O(n)$; 1, span $O(\log n)$; 2, why-span ---
pairwise combine up a $\log n$-height tree, levels parallel (``it's
parallel'' alone = 1).
\end{solution}

\textbf{(b)} \subpts{2} The function \code{reduce g z S} requires that
\code{g} be \textbf{associative}. A student claims \code{Seq.reduce (op
div) 1 S} computes the same result as folding \code{div} over the
corresponding list, for any \code{S}. Give a specific sequence \code{S}
for which this fails, and state the underlying reason in one sentence.

\answerspace{3.5cm}

\begin{solution}
\textbf{Counterexample:} take \code{S = }$\langle$\code{16, 4, 2}$\rangle$.
Folding \code{div} left-to-right groups as \code{(16 div 4) div 2}
$\eeq$ \code{4 div 2} $\eeq$ \code{2}, whereas \code{reduce} may
re-associate to the right as \code{16 div (4 div 2)} $\eeq$
\code{16 div 2} $\eeq$ \code{8}. Since \code{2} $\neq$ \code{8}, the two
disagree.

\textbf{Reason:} integer division isn't associative, so \code{reduce}'s
reparenthesization need not agree with a left-to-right fold.

\textbf{Rubric:} 1, a concrete \code{S} that diverges (no div-by-zero in
any grouping); 1, because \code{div} isn't associative.
\end{solution}

\newpage

% =============================================================================
% PROBLEM 4 --- Lazy Programming (eager-eval + maximal laziness recognition)
% =============================================================================

\quizproblem{Now You See It}{6}{6}

SML is \term{eagerly evaluated}: a function's arguments are fully
evaluated to values before the body runs. \term{Lazy} evaluation instead
computes an expression only when its value is needed. A \term{thunk}
(or \term{suspension}) is a value of type \code{unit -> 'a} used to
delay a computation.

\textbf{(a)} \subpts{2} What is the result of evaluating the following?
Circle one, and justify in one sentence.

\begin{codeblock}
  (fn (x, y) => x) (1, 1 div 0)
\end{codeblock}

\vspace{0.3em}
\hspace{2em} evaluates to \code{1} \qquad raises an exception \qquad
loops forever

\answerspace{2.5cm}

\begin{solution}
\textbf{Raises an exception} (\code{Div}). SML is eager, so the argument
tuple \code{(1, 1 div 0)} is evaluated before the projection runs, and
\code{1 div 0} raises \code{Div}.

\textbf{Rubric:} 1, raises \code{Div}; 1, eager eval forces the whole
tuple before the projection.
\end{solution}

\textbf{(b)} \subpts{4} Here are two ways to encode an infinite sequence
of values. A \term{lazy list} stores its head directly; a \term{stream}
hides its \emph{whole} front behind a suspension:

\begin{codeblock}
  datatype 'a llist  = LNil | LCons of 'a * (unit -> 'a llist)

  datatype 'a stream = Stream of (unit -> 'a front)
       and 'a front  = Empty  | Cons  of 'a * 'a stream
\end{codeblock}

Each has a \code{tabulate} that builds the sequence
\code{f 0, f 1, f 2, ...}:

\begin{codeblock}
  fun ltabFrom f i = LCons (f i, fn () => ltabFrom f (i + 1))
  fun ltabulate f  = ltabFrom f 0

  fun stabFrom f i = Stream (fn () => Cons (f i, stabFrom f (i + 1)))
  fun stabulate f  = stabFrom f 0
\end{codeblock}

Now suppose we evaluate the two bindings below, where \code{f i = 1024
div i} (so \code{f 0} raises \code{Div}):

\begin{codeblock}
  val a = ltabulate (fn i => 1024 div i)   (* lazy list *)
  val b = stabulate (fn i => 1024 div i)   (* stream    *)
\end{codeblock}

Exactly one of these bindings raises an exception. \textbf{Which one, and
why does the other not?} Refer to what each of \code{ltabulate} and
\code{stabulate} actually computes at the moment its binding is
evaluated.

\answerspace{5.5cm}

\begin{solution}
\textbf{\code{a} raises \code{Div}; \code{b} does not.} The difference is
whether the head element is behind a suspension.

\begin{itemize}
  \item \code{ltabulate f} builds \code{LCons (f 0, ...)}. The head
    \code{f 0} is a direct argument, so it is evaluated when \code{a} is
    bound, and \code{f 0 = 1024 div 0} raises \code{Div}. Only the tail
    is suspended.
  \item \code{stabulate f} builds \code{Stream (fn () => Cons (f 0,
    ...))}, with the head \code{f 0} inside the thunk. Binding \code{b}
    just stores that thunk (a value); \code{f 0} runs only when the
    stream is exposed, so no \code{Div} yet.
\end{itemize}

A lazy list forces its first element by construction; a stream defers
even that. That extra deferral is \term{maximal laziness}.

\textbf{Rubric:} 2, \code{a} raises / \code{b} doesn't; 1, \code{a}'s
head \code{f 0} is a direct arg to \code{LCons}, forced at bind time;
1, \code{b}'s head sits inside the thunk, forced only on expose.
\end{solution}

\end{document}
